\documentclass[11pt]{article}

\usepackage[margin=1in]{geometry}
\usepackage{amsmath,amssymb,amsthm}
\usepackage{microtype}
\usepackage{enumitem}
\usepackage[hidelinks]{hyperref}
\usepackage{xurl}
\usepackage{booktabs}
\usepackage{tabularx}

\title{Deployed Observation Surfaces for Anonymous DHTs\\
\large Control-plane knobs in IPFS/libp2p delegated routing and OHTTP evolution}
\author{}
\date{Unified archive note v0.20 (2026-03-06; archive v164)}

\begin{document}
\maketitle

\begin{abstract}
Anonymous-DHT papers often focus on protocol mechanisms (delegation, PIR, cover, padding).
Real deployments frequently leak through \emph{control-plane} and \emph{service-concentration} surfaces: shared delegated-routing endpoints, cache layers, endpoint autoconfiguration, feature negotiation, and fallback.
This note isolates those deployed surfaces for IPFS/libp2p and provides a compact ``what-to-pin'' checklist, so other papers in this archive can stay mechanism-focused and non-redundant.
\end{abstract}

\paragraph{Scope.}
We focus on IPFS/libp2p content/peer routing where ``anonymous lookup'' is pursued via delegated routing, role separation, and incremental privacy upgrades.
We treat these as \emph{observation surfaces} that must be declared (as projections, state declarations, or separation manifests) before any anonymity budget is meaningful.

\paragraph{Imports.}
Observation-tier naming and exposure normal forms (ENF-ID) are imported from Synthesis~3; threat-window semantics (TW-ID) from Synthesis~4.
Change-control equivalence for plan/state drift is imported from Synthesis~18.
Primary-path role separation is imported from Synthesis~22--23, and profiler/equalization contracts from Anonymity~B.
\cite{series:traceifaces,series:threatwindows,series:planstateequiv,series:primarysepmanifest,series:probsep,series:profiling}

\paragraph{Exports.}
Exports a compact deployment-surface taxonomy and a minimal pinning checklist: what belongs in \texttt{state\_decl\_id}, what belongs in a separation manifest, and what must be treated as part of the declared transcript projection $X=\pi(T)$.
\emph{Rule of thumb:} long-lived deployment state (router sets, cache layers, AutoConf refresh, downgrade rules) belongs in \texttt{state\_decl\_id}; a specific external service interface (method + URL + request-shaping defaults + time budgets) should be bound per line item via \texttt{contact\_surface\_id} (Synthesis~12).\cite{series:receiptitems}


\paragraph{How other papers should cite this note.}
This note owns the \emph{enumeration} of concrete deployed control-plane surfaces.
It does \emph{not} re-define receipts or state declarations: line-item tuple structure is standardized by Synthesis~12, and digest-bound change control for plan/state ids is standardized by Synthesis~18.\cite{series:receiptitems,series:planstateequiv}
Other papers should therefore cite this note only when they need to justify that a particular knob/end\-point/cache is part of the reader set (and hence must be pinned), and otherwise cite Synthesis~12/18 for the mechanics of pinning.

\paragraph{Three-way ownership rule (the compact version).}
When a deployed service is involved, this archive now uses a strict split:
(1) \texttt{state\_decl\_id} owns \emph{who may be contacted and how the endpoint set is resolved} (AutoConf, overrides, allow/deny lists, ignore-provider policies, timeout/fallback rules);
(2) \texttt{contact\_surface\_id} owns the named service interface actually invoked (method + full path-scoped URL + declared request-shaping defaults + declared time budgets);
(3) \texttt{exposure\_nf\_id} owns the \emph{reader-visible transcript semantics} (status-code interpretation, empty-result equivalence, streaming/chunking shape, and other projection details).
The worked example now carries a tiny adjunct making this split explicit so replay code does not have to reconstruct it from prose.

\section{Two kinds of deployed leakage}
This archive treats ``deployed'' leakage as falling into two dominant classes.

\subsection{Service concentration (shared readers)}
A mechanism may offer meaningful query privacy only if the \emph{reader} does not trivially see the initiator.
In practice, shared public endpoints and cache-backed services can concentrate observation at a small set of operators.
For example, IPFS Docs list a public delegated-routing endpoint (\texttt{https://delegated-ipfs.dev/routing/v1}) backed by \emph{someguy}; when used as a default, that single service becomes a reader unless an additional privacy layer is deployed.\cite{ipfsDocsPublicUtilities,ipfsSomeguy}
The same docs also describe \texttt{cid.contact} as a public Network Indexer that implements the same delegated-routing API, making it another shared reader surface.\cite{ipfsDocsNodes}
The IPNI docs further note that the public endpoint \texttt{https://delegated-ipfs.dev/routing/v1} is operated as a caching proxy for both Amino DHT and IPNI queries, forwarding to infrastructure such as \texttt{cid.contact}; this is exactly the kind of ``under the hood coupling that should be treated as a reader surface unless explicitly anonymized.\cite{ipfsDocsIPNI,ipfsSomeguy}
For anonymous-DHT accounting, the correct response is not prose about ``using delegated routing''; it is to (i) declare which service(s) are in the path and (ii) decide whether the design relies on \emph{separation} (non-collusion) or on a cryptographic/privacy layer that tolerates a single reader.

\subsection{Control-plane drift (silent surface changes)}
Even with a fixed mechanism, deployments drift.
If endpoint selection, feature negotiation, or fallback behavior changes over time, then the observable transcript projection changes.
This archive's rule is: drift changes claim identity.
Bind such knobs into a digest-bound plan/state declaration (Synthesis~18) so changes are diffable by id, not discovered by inference.
\cite{series:planstateequiv}

\paragraph{Concrete drift example (public indexer incident response).}
In September 2025, \texttt{cid.contact} announced it would add a temporary set of delegated-routing endpoints to restore publisher reachability during an IPNI disruption.
This is a clean example of why endpoint-set drift must be digest-bound and reviewable: the reader set and the caching path can change quickly for operational reasons, and that change is privacy-relevant.\cite{ipfsDiscussCidContactEndpoints2025}

\section{Delegated routing (Routing V1) as a surface}
Delegated routing is a vendor-agnostic HTTP API used to offload content routing / peer routing / naming to a service.\cite{ipfsDelegatedRoutingV1Spec,ipip0337}
Deployed privacy questions usually reduce to three concrete surfaces:

\begin{enumerate}[leftmargin=*]
\item \textbf{Endpoint identity.} Which delegated-routing base URLs are used (including localhost/self-hosted routing backends)? Are they public shared endpoints? Are they pinned or discoverable?
IPFS Docs explicitly names a public-good delegated-routing endpoint at \texttt{https://delegated-ipfs.dev/routing/v1} (operated by the IPFS Foundation and backed by \emph{someguy}) and notes that \texttt{cid.contact} also implements the Delegated Routing V1 API.\cite{ipfsDocsNodes}
Routing~V1 also supports \emph{path-scoped capability endpoints} such as \texttt{/routing/v1/providers}, \texttt{/routing/v1/peers}, and \texttt{/routing/v1/ipns} (with the no-path form enabling all operations). Kubo documentation shows AutoConf manifests that select different endpoint \emph{paths} per capability and even per node class (``with DHT'' vs ``without DHT''). Changing a path changes the observable transcript surface and must be treated as semantic drift of the contact surface, not as a cosmetic URL change.\cite{kuboConfigPathRouting,ipfsDelegatedRoutingV1Spec}
Since Kubo v0.40.0 (and still in the v0.40.1 patch line), the Routing~V1 HTTP API is exposed by default on the gateway port at \texttt{http://127.0.0.1:8080/routing/v1}, allowing light clients (e.g., browsers) to treat a local gateway as a delegated-routing backend.\cite{kuboReleaseV040,kuboReleaseV0401}
This can be disabled via \texttt{Gateway.ExposeRoutingAPI}. In a public gateway deployment, leaving it enabled changes the reader set (who sees routing queries) and therefore must be pinned as part of \texttt{state\_decl\_id}.\cite{kuboConfigExposeRoutingAPI}
\item \textbf{Cache layers.} Are requests served from a shared cache (CDN / edge / reverse proxy) that sees query patterns? If yes, the cache operator becomes an observation point.
The IPFS blog describes caching in the public delegated-routing endpoint to speed repeated browser lookups; this also concentrates observation at the cache operator.\cite{blogDelegatedRoutingCaching2025}
In addition, the Helia delegated-routing client caches successful (HTTP~200) responses in browser environments (Cache API) for a short TTL, creating a second, often-overlooked cache layer that affects what is observable when.\cite{heliaDelegatedRoutingClientCache}
\item \textbf{Feature growth.} The Routing~V1 family evolves (streaming responses, client filtering, closest-peers endpoints).\cite{ipip0410,ipip0484,ipip0476}
Each new endpoint potentially enlarges the declared projection $X$ (message count, response shape, pacing).
Even without a new endpoint, response semantics can drift: the current Routing~V1 interoperability guidance moves empty results toward HTTP~200 with an empty collection, while requiring clients to interpret legacy HTTP~404 as an empty-result equivalent. That is best modeled as projection/ENF drift, not merely a hidden state change, because the reader-visible transcript shape changed even if the URL stayed fixed.\cite{ipip0513,ipfsDelegatedRoutingV1Spec}
\end{enumerate}

\paragraph{Pinning rule.}
The base URL set and any ``auto'' resolution rules belong in \texttt{state\_decl\_id}.
If the privacy story assumes a relay/service split, encode it as a separation manifest (Synthesis~22) rather than a prose promise.
\cite{series:primarysepmanifest}

\section{AutoConf and ``auto'' placeholders}
Kubo supports \emph{AutoConf}: fetching a remote JSON manifest (\texttt{autoconf.json}) that supplies network defaults (bootstrap peers, delegated routers, DNS resolvers, etc.) and resolves explicit \texttt{["auto"]} placeholders at runtime.\cite{kuboConfigAutoconf,kuboChangelogV037,kuboRPCexpandAuto,ipfsMainnetAutoConfJSON,ipfsMainnetAutoConfSite}
AutoConf is not only a convenience feature: it is a control-plane dependency that can silently change which delegated-routing endpoints are used (and therefore who reads queries). On IPFS Mainnet, the default manifest is machine-readable and includes delegated endpoints such as \texttt{https://delegated-ipfs.dev} and \texttt{https://cid.contact}.\cite{ipfsMainnetAutoConfJSON}

\paragraph{Override precedence (routing endpoint selection).}
For receipt purposes, treat \emph{the effective router set} as the outcome of a small precedence stack:
\begin{center}
\small
\begin{tabularx}{0.98\linewidth}{@{}lX@{}}
\toprule
Precedence (high$\to$low) & What it can change (and why it must be pinned) \\
\midrule
Process environment (e.g., \texttt{IPFS\_HTTP\_ROUTERS}) & Replaces delegated-router configuration and AutoConf expansion; shifts the reader set immediately.\\
Explicit config (router backends + per-capability paths) & Chooses concrete router implementations (DHT vs HTTP routing) and path-scoped endpoints; changes transcript shape and termination.\\
AutoConf \texttt{auto} expansion & Resolves \texttt{["auto"]} placeholders via remote manifest; can silently rotate endpoints over time.\\
\bottomrule
\end{tabularx}
\end{center}
The rule for this series is simple: if a knob participates in this stack, it belongs in \texttt{state\_decl\_id}.\cite{series:planstateequiv,series:receiptitems}


Kubo also exposes a hard override: \texttt{IPFS\_HTTP\_ROUTERS} replaces AutoConf and any manually configured delegated routers when set.\cite{kuboEnvVars}

Kubo's environment-variable layer also exposes a request-shaping override: \texttt{IPFS\_HTTP\_ROUTERS\_\allowbreak FILTER\_\allowbreak PROTOCOLS} overrides the values sent via the \texttt{filter-protocols} parameter defined by IPIP-0484 for Delegated Routing V1 requests. This changes which transfer protocols are returned (and therefore which retrieval path is feasible) and must be treated as declared control-plane state.
\cite{kuboEnvVars,ipip0484,ipfsDelegatedRoutingV1Spec}

Other delegated-routing clients may implement IPIP-0484 protocol filtering either as a fixed default or as an explicit knob.
If a deployed client hardcodes a filter set, that set is effectively part of the binary and must be pinned by release/commit inside \texttt{state\_decl\_id}.
Rainbow's documented configuration surface (as of 2026-03-05) includes router endpoint selection and time budgets but does not document a protocol-filter knob; deployments that rely on protocol filtering should therefore pin the deployed client version or explicitly declare whether \texttt{filter-protocols} is applied.
\cite{rainbowEnvVars,ipip0484}

Rainbow exposes the rest of this routing surface as first-class configuration: \texttt{RAINBOW\_HTTP\_ROUTERS} selects the delegated-routing endpoint set, \texttt{RAINBOW\_AUTOCONF\_URL} selects the manifest, \texttt{RAINBOW\_AUTOCONF\_\allowbreak REFRESH} controls the rotation cadence for \texttt{auto} placeholders, \texttt{RAINBOW\_HTTP\_ROUTERS\_\allowbreak TIMEOUT} controls the network-level timeout for individual router HTTP requests, and \texttt{RAINBOW\_ROUTING\_\allowbreak TIMEOUT} controls the application-level time budget for a parallel routing operation. These knobs are privacy-relevant because they can change which services are contacted, how long a primary attempt is allowed to run, and therefore whether (and how often) fallback triggers; they should be pinned (or explicitly excluded) in \texttt{state\_decl\_id}.\cite{rainbowEnvVars}

At the process level, \texttt{Routing.Routers} selects the concrete router backends (HTTP delegated routers, DHT, etc.), and \texttt{Routing.AcceleratedDHTClient} changes the contact pattern and termination behavior of the DHT client.\cite{kuboConfigPathRouting}

\paragraph{Concrete drift example (defaults move to \texttt{auto}).}
The Rainbow gateway implementation documents that its default delegated-routing setting moved from a fixed endpoint (\texttt{https://cid.contact}) to \texttt{auto} expansion via AutoConf.
This illustrates how a seemingly ``minor'' config default change can move a deployment from a single named reader to a dynamically selected endpoint set.\cite{rainbowChangelogAutoconf}

\paragraph{Self-hosted gateways as delegated-routing backends.}
If a deployment uses Kubo gateways as delegated-routing backends (the v0.40.0 default), the configuration knob \texttt{Gateway.ExposeRoutingAPI} and the gateway's access-control posture belong in the declared state surface alongside any AutoConf resolution rules; otherwise the reader set can change without a claim-id change.\cite{kuboReleaseV040,kuboConfigExposeRoutingAPI}

From an anonymity perspective, AutoConf is a \emph{control-plane dependency}:
\begin{itemize}[leftmargin=*]
\item it can silently change which endpoints you talk to (surface drift),
\item it can concentrate observation if defaults point to a small set of public endpoints,
\item it can couple privacy posture to availability events (fallback/downgrade).
\end{itemize}

\paragraph{Pinning rule (AutoConf).}
If AutoConf influences the routing surface, then (at minimum) the AutoConf URL, the resolution rules for \texttt{["auto"]}, the refresh cadence (e.g., Rainbow \texttt{RAINBOW\_AUTOCONF\_\allowbreak REFRESH}), any \emph{override knobs} (e.g., \texttt{IPFS\_HTTP\_ROUTERS}), and the allowed fallback behavior belong in \texttt{state\_decl\_id}.
If the deployment claims ``primary-path'' behavior (e.g., privacy relays enabled) then downgrade/fallback is a first-class leakage event and must be pinned by the schedule/state discipline, not left implicit.\cite{series:operationalA,series:planstateequiv}

\section{HTTP retrieval is a derived routing surface}
Delegated-routing lookups do not just reveal \emph{where} queries go; they can also change \emph{how} blocks are retrieved.
Kubo exposes \texttt{HTTPRetrieval.Enabled}: when enabled, it will act on \texttt{/tls/http} provider multiaddrs returned by \texttt{Routing.DelegatedRouters} and perform pure HTTPS block retrieval (trustless-gateway style \verb|GET /ipfs/{cid}| with \texttt{format=raw} or \texttt{Accept: application/vnd.ipld.raw}), alongside or instead of Bitswap.\cite{kuboConfigHTTPRetrieval,trustlessGatewaySpec}
This upgrades the ``destination'' surface from ``libp2p peers'' to ``HTTP services'' (often with their own caches, CDNs, and logs), and should be pinned as state whenever destination privacy is claimed.
Rainbow exposes an analogous destination surface: \texttt{RAINBOW\_HTTP\_RETRIEVAL\_ENABLE} gates HTTP retrieval, and allow/deny lists (e.g., \texttt{RAINBOW\_HTTP\_RETRIEVAL\_ALLOWLIST}, \texttt{RAINBOW\_HTTP\_RETRIEVAL\_DENYLIST}) determine which HTTP endpoints are eligible; timeouts (\texttt{RAINBOW\_RETRIEVAL\_TIMEOUT}) influence termination and downgrade behavior. Treat these as declared state whenever the receipt's destination model assumes a particular retrieval regime.\cite{rainbowEnvVars}
Kubo makes the same boundary explicit on the config side: \texttt{HTTPRetrieval.Allowlist} and \texttt{HTTPRetrieval.Denylist} filter destination hostnames, while \texttt{Routing.IgnoreProviders} suppresses specific provider peer IDs. Rainbow exposes the peer-id side as \texttt{ROUTING\_IGNORE\_PROVIDERS}. These are not cosmetic tuning knobs: they change which destinations are ever eligible to observe retrieval traffic and therefore belong under change control.\cite{kuboConfigHTTPRetrieval,kuboConfigPathRouting,rainbowEnvVars}


For Anonymous-DHT receipts, it is usually enough to declare (i) whether HTTP retrieval is enabled (Kubo \texttt{HTTPRetrieval.Enabled}; Rainbow \texttt{RAINBOW\_HTTP\_RETRIEVAL\_ENABLE}), (ii) the accepted protocol/filter rules (including any \texttt{filter-protocols} override)\cite{ipip0484,kuboEnvVars}, (iii) destination allow/deny and ignore-provider policies, and (iv) retrieval timeout/abort semantics (Kubo retrieval timeouts; Rainbow \texttt{RAINBOW\_RETRIEVAL\_TIMEOUT}) as part of the threat window.\cite{series:threatwindows,rainbowEnvVars,kuboConfigHTTPRetrieval,kuboConfigPathRouting}

\section{OHTTP upgrades and chunked-message evolution}
Reader privacy for delegated routing contemplates Oblivious HTTP relays.\cite{rfc9458OHTTP,ipip0421Draft}
Discovery and configuration can be standardized (OHTTP discovery).\cite{rfc9540OHTTPDiscovery}
The OHTTP application layer is also evolving: \emph{chunked} OHTTP message formats enable incremental processing.\cite{draftChunkedOHTTP}

\paragraph{Why chunking matters.}
Chunking introduces new observable features (chunk counts, boundaries, pacing, partial-response timing) that can become profiling handles unless they are explicitly treated as part of the declared projection and budgeted (Synthesis~3; Anonymity~B).\cite{series:traceifaces,series:profiling}

\section{Minimal checklist (what-to-pin)}
This is the ``copy/paste'' checklist for other papers.

\begin{table}[t]
\centering
\begin{tabularx}{\linewidth}{@{}>{\raggedright\arraybackslash}X >{\raggedright\arraybackslash}X@{}}
\toprule
Surface & Where it must appear in the published claim interface \\
\midrule
Delegated-routing endpoint set (full URLs incl. path scopes) & \texttt{state\_decl\_id} (digest-bound); drift forces a new claim id. \\
Router set + overrides (Kubo \texttt{Routing.DelegatedRouters}/AutoConf/\texttt{IPFS\_HTTP\_ROUTERS}; Rainbow \texttt{RAINBOW\_HTTP\_ROUTERS} + \texttt{RAINBOW\_AUTOCONF\_\allowbreak REFRESH} + \texttt{RAINBOW\_HTTP\_ROUTERS\_\allowbreak TIMEOUT} + \texttt{RAINBOW\_ROUTING\_\allowbreak TIMEOUT}) & \texttt{state\_decl\_id}; drift changes the reader set and the transcript. \\
AutoConf URL and \texttt{auto} resolution rules & \texttt{state\_decl\_id}; include fallback/downgrade semantics. \\
HTTP retrieval + protocol filters (Kubo \texttt{HTTPRetrieval.Enabled}, \texttt{HTTPRetrieval.Allowlist}/\texttt{Denylist}, \texttt{Routing.IgnoreProviders}, \texttt{filter-protocols}, \texttt{IPFS\_HTTP\_ROUTERS\_\allowbreak FILTER\_\allowbreak PROTOCOLS}; Rainbow \texttt{RAINBOW\_HTTP\_RETRIEVAL\_ENABLE}, allow/deny lists, \texttt{ROUTING\_IGNORE\_PROVIDERS}, and pinned client version for any implicit protocol-filter defaults) & \texttt{state\_decl\_id} or explicit line item; drift changes the destination reader set. \\
Cache/CDN/reverse-proxy readers & If relied upon: separation manifest; otherwise treat as a single reader in the threat model and budget accordingly. \\
Feature negotiation (streaming/filtering/closest-peers endpoints, empty-result status semantics) & Treat as part of the declared projection $X$ (ENF-ID) and under change control. \\
OHTTP relay split (client/relay/target non-collusion) & Separation manifest (Synthesis~22) and (if probabilistic) $(0,\rho)$ budget (Synthesis~23). \\
OHTTP chunking mode and its pacing/shape affordances & Declared projection $X$ + pinned plan/state; otherwise it is uncontrolled feature drift. \\
\bottomrule
\end{tabularx}
\caption{Deployed surfaces and where to pin them in the One-True-Archive interface stack.}
\label{tab:pinning}
\end{table}

\section{Relation to the related-work map}
This note is intentionally not a literature taxonomy.
For representative academic comparators (delegation, PIR, obfuscation, measurement), cite the related-work map (Synthesis~7).\cite{series:relatedworkmap}

\begin{thebibliography}{99}\small

\bibitem{series:traceifaces}
Anonymous.
\newblock Trace Interfaces for Anonymous DHTs: Observation Tiers, Committee-Contact Traces, and Exposure Normal Forms.
\newblock Series synthesis note (Synthesis~3), 2026. (This archive.)

\bibitem{series:threatwindows}
Anonymous.
\newblock Threat Windows for Anonymous DHT Lookups: Horizon Semantics, Retry Surfaces, and Declared Templates.
\newblock Series synthesis note (Synthesis~4), 2026. (This archive.)

\bibitem{series:planstateequiv}
Anonymous.
\newblock Digest-Bound Replay Plans and State Declarations: Equivalence, Minimality, and Change Control.
\newblock Series synthesis note (Synthesis~18), 2026. (This archive.)

\bibitem{series:receiptitems}
Anonymous.
\newblock Receipt Line Items for Anonymity Claims: A Minimal Tuple Schema for Published Budgets.
\newblock Series synthesis note (Synthesis~12), 2026. (This archive.)

\bibitem{series:primarysepmanifest}
Anonymous.
\newblock Primary-Path Separation Manifests: Digest-Bound Non-Collusion Assumptions for Anonymous DHT Claims.
\newblock Series synthesis note (Synthesis~22), 2026. (This archive.)

\bibitem{series:probsep}
Anonymous.
\newblock Probabilistic Separation Budgets: $(0,\rho)$ Checkpoints for Non-Collusion Drift.
\newblock Series synthesis note (Synthesis~23), 2026. (This archive.)

\bibitem{series:profiling}
Anonymous.
\newblock Anonymous DHT Profiling and the Equalization Contract.
\newblock Anonymity series Paper~B, 2026. (This archive.)

\bibitem{series:operationalA}
Anonymous.
\newblock Schedule-Only Retries for Anonymous DHT Lookups.
\newblock Operational series Paper~A, 2026. (This archive.)

\bibitem{series:relatedworkmap}
Anonymous.
\newblock Related Work Map for Anonymous DHT Lookups.
\newblock Series synthesis note (Synthesis~7), 2026. (This archive.)

\bibitem{ipfsDelegatedRoutingV1Spec}
G.~Eggert, M.~H. Derkani, H.~Dias, and M.~Rataj.
\newblock \emph{Delegated Routing V1 HTTP API}.
\newblock IPFS Standards, 17 Dec 2025.
\newblock \url{https://specs.ipfs.tech/routing/http-routing-v1/}.

\bibitem{ipfsDocsNodes}
IPFS Docs.
\newblock \emph{Nodes: public delegated routing endpoints and indexers}.
\newblock Documentation, 2025--2026.
\newblock \url{https://docs.ipfs.tech/concepts/nodes/}.

\bibitem{ipfsDocsPublicUtilities}
IPFS Docs.
\newblock \emph{Public IPFS Utilities}.
\newblock Documentation, 2025--2026.
\newblock \url{https://docs.ipfs.tech/concepts/public-utilities/}.

\bibitem{ipfsDocsIPNI}
IPFS Docs.
\newblock \emph{IPNI (InterPlanetary Network Indexer)}.
\newblock Documentation, 2025--2026.
\newblock \url{https://docs.ipfs.tech/concepts/ipni/}.

\bibitem{ipfsSomeguy}
IPFS Project.
\newblock \emph{someguy: Delegated Routing V1 server (Amino DHT + Network Indexer proxy)}.
\newblock Source repository, 2023--2026.
\newblock \url{https://github.com/ipfs/someguy}.

\bibitem{blogDelegatedRoutingCaching2025}
IPFS Blog.
\newblock \emph{Faster Peer-to-Peer Retrieval in Browsers With Caching in the Public Delegated Routing Endpoint}.
\newblock 5 Sep 2025.
\newblock \url{https://blog.ipfs.tech/2025-delegated-routing-caching/}.

\bibitem{ipfsDiscussCidContactEndpoints2025}
IPFS Discuss.
\newblock \emph{Adding delegated routing endpoints to cid.contact}.
\newblock 10 Sep 2025.
\newblock \url{https://discuss.ipfs.tech/t/adding-delegated-routing-endpoints-to-cid-contact/19727}.


\bibitem{ipip0337}
IPIP-0337 Authors.
\newblock \emph{IPIP-0337: Delegated Content Routing HTTP API}.
\newblock IPFS Standards (IPIP), 18 Oct 2022.
\newblock \url{https://specs.ipfs.tech/ipips/ipip-0337/}.


\bibitem{ipip0410}
IPIP-0410 Authors.
\newblock \emph{IPIP-0410: Streaming NDJSON in Routing HTTP API}.
\newblock IPFS Standards (IPIP), 12 May 2023.
\newblock \url{https://specs.ipfs.tech/ipips/ipip-0410/}.


\bibitem{ipip0484}
IPIP-0484 Authors.
\newblock \emph{IPIP-0484: Opt-in Filtering in Routing V1 HTTP API}.
\newblock IPFS Standards (IPIP), 29 Oct 2024.
\newblock \url{https://specs.ipfs.tech/ipips/ipip-0484/}.


\bibitem{ipip0476}
IPIP-0476 Authors.
\newblock \emph{IPIP-0476: Delegated Routing DHT Closest Peers API}.
\newblock IPFS Standards (IPIP), 20 Nov 2025.
\newblock \url{https://specs.ipfs.tech/ipips/ipip-0476/}.


\bibitem{ipip0513}
IPIP-0513 Authors.
\newblock \emph{IPIP-0513: Empty Results Response Update for Routing V1 HTTP API}.
\newblock IPFS Standards (IPIP), 13 Feb 2026.
\newblock \url{https://specs.ipfs.tech/ipips/ipip-0513/}.


\bibitem{ipfsMainnetAutoConfJSON}
IPFS Mainnet AutoConf.
\newblock \emph{Canonical AutoConf manifest (autoconf.json)}.
\newblock Accessed 2026-02-28.
\newblock \url{https://conf.ipfs-mainnet.org/autoconf.json}.

\bibitem{kuboConfigAutoconf}
IPFS Kubo.
\newblock \emph{Kubo configuration: AutoConf}.
\newblock Documentation (raw), accessed 2026-03-05.
\newblock \url{https://raw.githubusercontent.com/ipfs/kubo/master/docs/config.md}.


\bibitem{kuboRPCexpandAuto}
IPFS Docs.
\newblock \emph{Kubo RPC API v0 reference: \texttt{expand-auto} resolves \texttt{auto} placeholders from AutoConf}.
\newblock Updated 2025; accessed 2026-02-28.
\newblock \url{https://docs.ipfs.tech/reference/kubo/rpc/}.


\bibitem{kuboEnvVars}
IPFS Kubo.
\newblock \emph{Kubo environment variables: \texttt{IPFS\_HTTP\_ROUTERS}, \texttt{IPFS\_HTTP\_ROUTERS\_\allowbreak FILTER\_\allowbreak PROTOCOLS}, and routing overrides}.
\newblock Documentation (raw), accessed 2026-03-05.
\newblock \url{https://raw.githubusercontent.com/ipfs/kubo/master/docs/environment-variables.md}.


\bibitem{kuboChangelogV037}
IPFS Kubo.
\newblock \emph{Kubo v0.37 changelog: AutoConf control and explicit \texttt{["auto"]} placeholders}.
\newblock 27 Aug 2025.
\newblock \url{https://raw.githubusercontent.com/ipfs/kubo/master/docs/changelogs/v0.37.md}.

\bibitem{rainbowChangelogAutoconf}
IPFS Rainbow.
\newblock \emph{CHANGELOG: AutoConf default \texttt{auto} expansion and control-plane default changes}.
\newblock 2025--2026; accessed 2026-02-28.
\newblock \url{https://github.com/ipfs/rainbow/blob/main/CHANGELOG.md}.


\bibitem{ipfsMainnetAutoConfSite}
IPFS Mainnet AutoConf.
\newblock \emph{Registry and format notes (alpha; subject to change)}.
\newblock 2025--2026; accessed 2026-02-28.
\newblock \url{https://conf.ipfs-mainnet.org/}.


\bibitem{rfc9458OHTTP}
T.~Pauly, M.~Thomson, and C.~A. Wood.
\newblock Oblivious HTTP.
\newblock RFC~9458, IETF, 2023.

\bibitem{rfc9540OHTTPDiscovery}
M.~Thomson and C.~A. Wood.
\newblock Discovering Oblivious Services via HTTP.
\newblock RFC~9540, IETF, 2024.

\bibitem{ipip0421Draft}
I.~Schasny and contributors.
\newblock \emph{Draft IPIP-0421: HTTP Delegated Routing Reader Privacy Upgrade}.
\newblock Pull request \#421 in \emph{ipfs/specs}, June 2023 (draft / changes requested).
\newblock \url{https://github.com/ipfs/specs/pull/421}.


\bibitem{heliaDelegatedRoutingClientCache}
IPFS Helia.
\newblock \emph{@helia/delegated-routing-v1-http-api-client documentation (browser caching defaults)}.
\newblock API docs, accessed 2026-03-04.
\newblock \url{https://ipfs.github.io/helia-delegated-routing-v1-http-api/modules/_helia_delegated-routing-v1-http-api-client.html}.

\bibitem{draftChunkedOHTTP}
D.~Schinazi, C.~Wood, and J.~Thomson (editors).
\newblock \emph{Chunked Oblivious HTTP Messages} (Internet-Draft).
\newblock IETF, draft-ietf-ohai-chunked-ohttp (latest), 2025--2026.
\newblock \url{https://datatracker.ietf.org/doc/draft-ietf-ohai-chunked-ohttp/}.

\bibitem{kuboReleaseV0401}
IPFS Kubo.
\newblock \emph{Release v0.40.1}.
\newblock GitHub releases page (patch line), 27 Feb 2026.
\newblock \url{https://github.com/ipfs/kubo/releases/tag/v0.40.1}.

\bibitem{kuboReleaseV040}
IPFS Kubo.
\newblock \emph{Release v0.40.0: gateway exposes Routing~V1 at \texttt{/routing/v1} by default}.
\newblock 25 Feb 2026.
\newblock \url{https://github.com/ipfs/kubo/releases/tag/v0.40.0}.

\bibitem{kuboConfigExposeRoutingAPI}
IPFS Kubo.
\newblock \emph{Kubo configuration: \texttt{Gateway.ExposeRoutingAPI}}.
\newblock Documentation (raw), accessed 2026-02-28.
\newblock \url{https://raw.githubusercontent.com/ipfs/kubo/master/docs/config.md}.

\bibitem{kuboConfigPathRouting}
IPFS Kubo.
\newblock Kubo configuration reference (\texttt{Routing} section: path-based delegated routing, AutoConf structure, and accelerated DHT client caveats).
\newblock \url{https://github.com/ipfs/kubo/blob/master/docs/config.md} (accessed 2026-03-05).



\bibitem{rainbowEnvVars}
IPFS Rainbow.
\newblock \emph{Rainbow environment variables: delegated routing, \texttt{RAINBOW\_HTTP\_ROUTERS}, and autoconf}.
\newblock Documentation, accessed 2026-03-05.
\newblock \url{https://github.com/ipfs/rainbow/blob/main/docs/environment-variables.md}.

\bibitem{kuboConfigHTTPRetrieval}
IPFS Kubo.
\newblock Kubo configuration reference (\texttt{HTTPRetrieval.Enabled}).
\newblock \url{https://github.com/ipfs/kubo/blob/master/docs/config.md} (accessed 2026-03-05).

\bibitem{trustlessGatewaySpec}
IPFS Standards.
\newblock \emph{Trustless Gateway Specification}.
\newblock \url{https://specs.ipfs.tech/http-gateways/trustless-gateway/} (accessed 2026-03-05).

\end{thebibliography}

\end{document}
